% GENERATED FILE. Edit canonical lessons, then run npm run generate:books.
% canonical-source-hash: fnv1a64:2c60faa7ecdc9d30
% canonical-lessons: ZH-C128-yinyue, ZH-W128-yin-sound, ZH-C128-ge, ZH-C128-haoting, ZH-C128-yinyuehui, ZH-C128-jita, ZH-R128-review-123-128

\chapter{Music}
\label{ch:zh-music}
\hlchaptermodality{128}

\begin{chapteropening}
\textbf{By the end of this chapter:} \emph{I can talk about music: songs, the piano and the guitar, and what sounds good.}

Complete the last lesson of chapter 128: I can talk about music: songs, the piano and the guitar, and what sounds good.
\end{chapteropening}

\section[\texorpdfstring{\zh{音乐} (yīnyuè)}{音乐 (yīnyuè)}]{\zh{音乐} (yīnyuè) — music}

\label{lesson:ZH-C128-yinyue}



\begin{quote}
Before the new one: say the Mandarin for a few, then the Mandarin for like this.
\end{quote}

\subsection*{You'll want to know: \zh{音乐}}
\textbf{\zh{音乐}} — \emph{yīnyuè} — \textquotedblleft{}music\textquotedblright{}.

\begin{quote}
\textbf{\zh{我喜欢音乐}} \emph{wǒ xǐhuan yīnyuè} — I love music.
\end{quote}

The new character, in pieces: \zh{音} — on top a five-stroke piece — a dot, a short bar, two short strokes leaning toward each other, and a long bar — and under it \zh{日}. The next lesson puts it in your hand.

\begin{cousinweb}
\textbf{\zh{音乐}} (\emph{yīnyuè}) is built from \zh{音} 'sound' and \zh{乐}, here said \emph{yuè}, both inherited from Middle Chinese, \emph{'im} and \emph{ngaewk} (Guangyun, in Baxter's transcription); the dictionary glosses the second with the whole word, 'music'.
\end{cousinweb}

\begin{grammarlens}[title={What you've built}]
One more word for what I like.
\end{grammarlens}

\subsection*{Your turn}
\begin{itemize}
\raggedright
  \item \emph{Say it:} \emph{yīnyuè}
  \item \emph{Say it:} \emph{yīnyuè}, once more
  \item \emph{Say it:} \emph{wǒ xǐhuan yīnyuè} — I love music.
  \item \emph{Recall:} say the Mandarin for a few, then the Mandarin for like this, then say \emph{yīnyuè} again
\end{itemize}

\begin{tcolorbox}[breakable,colback=teal!4,colframe=teal!35!black,title={Before you move on}]
What does \zh{音乐} mean? (\textquotedblleft{}Music\textquotedblright{}.) Say it once more.
\end{tcolorbox}
\section[\texorpdfstring{\zh{音} (yīn)}{音 (yīn)}]{\zh{音} — five strokes over \zh{日}}

\label{lesson:ZH-W128-yin-sound}



\begin{quote}
The word first, then the character inside it.

\begin{itemize}
\raggedright
  \item \emph{Recall:} say \emph{yīnyuè} — \textbf{R1}, one lesson back, before it decays
\end{itemize}

\begin{quote}
\zh{日} \zh{乐}
\end{quote}

\emph{Read it:} each one aloud, with nothing said first — \emph{rì}, \emph{lè}
\end{quote}

\subsection*{Script — five strokes over \zh{日}}
\hlblockfigure{\detokenize{figures/ZH-W128-yin-sound-filmstrip.pdf}}{How \zh{音} is written, stroke by stroke}

\begin{quote}
\zh{音}
\end{quote}

\textbf{Two pieces, top to bottom.}

\textbf{The top piece, five strokes.} A dot at the top, a short horizontal under it, a short stroke on the left falling to the lower right, a short stroke on the right falling to the lower left, then a long horizontal across the whole width.

\textbf{\zh{日}, four strokes}, the four your hand makes in \zh{日}: the left side, the top that turns down the right side, the middle bar, and the bar that closes the bottom.

5 + 4 = nine.

\subsection*{Your turn}
\begin{itemize}
\raggedright
  \item \emph{Write it:} \zh{音}, all nine, with the model in front of you
  \item \emph{Cover:} write it again from memory, and count the pieces, not the strokes
\end{itemize}

\begin{tcolorbox}[breakable,colback=teal!4,colframe=teal!35!black,title={Before you move on}]
How many pieces in \zh{音}? (\textbf{Two}.) Which stroke ends the top piece? (\textbf{The long horizontal}.) What sits underneath? (\textbf{\zh{日}}.)

Source: \href{https://raw.githubusercontent.com/chanind/hanzi-writer-data/68d10a4b21150cae5e1ebbd223eed289cf32d90c/data/\%E9\%9F\%B3.json}{Hanzi Writer Data \zh{音}, pinned stroke paths and medians}.
\end{tcolorbox}
\section[gē]{gē — song}

\label{lesson:ZH-C128-ge}



\begin{quote}
Before the new one: say the Mandarin for like this, then the Mandarin for music.
\end{quote}

\subsection*{You'll want to know: gē}
\emph{gē} — \textquotedblleft{}song\textquotedblright{}.

\begin{quote}
\emph{wǒ xǐhuan Zhōngguó gē} — I love Chinese songs.
\end{quote}

Its characters come later; for now it is a word for the ear.

\begin{cousinweb}
\emph{gē} is inherited from Middle Chinese \emph{ka} (Guangyun, in Baxter's transcription); the dictionary quotes an older book: the voice of a person is called song.
\end{cousinweb}

\begin{grammarlens}[title={What you've built}]
One more word for what I like.
\end{grammarlens}

\subsection*{Your turn}
\begin{itemize}
\raggedright
  \item \emph{Say it:} \emph{gē}
  \item \emph{Say it:} \emph{gē}, once more
  \item \emph{Say it:} \emph{wǒ xǐhuan Zhōngguó gē} — I love Chinese songs.
  \item \emph{Recall:} say the Mandarin for like this, then the Mandarin for music, then say \emph{gē} again
\end{itemize}

\begin{tcolorbox}[breakable,colback=teal!4,colframe=teal!35!black,title={Before you move on}]
What does gē mean? (\textquotedblleft{}Song\textquotedblright{}.) Say it once more.
\end{tcolorbox}
\section[hǎotīng]{hǎotīng — pleasant to hear}

\label{lesson:ZH-C128-haoting}



\begin{quote}
Before the new one: say the Mandarin for music, then the Mandarin for song.
\end{quote}

\subsection*{You'll want to know: hǎotīng}
\emph{hǎotīng} — \textquotedblleft{}pleasant to hear\textquotedblright{}.

\begin{quote}
\emph{Zhōngguó yīnyuè hěn hǎotīng} — Chinese music is lovely to listen to.
\end{quote}

Its characters come later; for now it is a word for the ear.

\begin{cousinweb}
\emph{hǎotīng} is built from \zh{好} 'good' and \emph{tīng} 'to listen', both inherited from Middle Chinese, \emph{xawX} and \emph{theng} (Guangyun, in Baxter's transcription), the same shape as \zh{好看} and \zh{好吃}.
\end{cousinweb}

\begin{grammarlens}[title={What you've built}]
One more word for what I like.
\end{grammarlens}

\subsection*{Your turn}
\begin{itemize}
\raggedright
  \item \emph{Say it:} \emph{hǎotīng}
  \item \emph{Say it:} \emph{hǎotīng}, once more
  \item \emph{Say it:} \emph{Zhōngguó yīnyuè hěn hǎotīng} — Chinese music is lovely to listen to.
  \item \emph{Recall:} say the Mandarin for music, then the Mandarin for song, then say \emph{hǎotīng} again
\end{itemize}

\begin{tcolorbox}[breakable,colback=teal!4,colframe=teal!35!black,title={Before you move on}]
What does hǎotīng mean? (\textquotedblleft{}Pleasant to hear\textquotedblright{}.) Say it once more.
\end{tcolorbox}
\section[\texorpdfstring{\zh{音乐会} (yīnyuèhuì)}{音乐会 (yīnyuèhuì)}]{\zh{音乐会} (yīnyuèhuì) — concert}

\label{lesson:ZH-C128-yinyuehui}



\begin{quote}
Before the new one: say the Mandarin for song, then the Mandarin for pleasant to hear.
\end{quote}

\subsection*{You'll want to know: \zh{音乐会}}
\textbf{\zh{音乐会}} — \emph{yīnyuèhuì} — \textquotedblleft{}concert\textquotedblright{}.

\begin{quote}
\textbf{\zh{明天晚上我们去音乐会}} \emph{míngtiān wǎnshang wǒmen qù yīnyuèhuì} — Tomorrow evening we're going to a concert.
\end{quote}

\begin{cousinweb}
\textbf{\zh{音乐会}} (\emph{yīnyuèhuì}) is built from \zh{音乐} 'music' and \zh{会} 'a gathering', all three parts inherited from Middle Chinese, \emph{'im}, \emph{ngaewk} and \emph{hwajH} (Guangyun, in Baxter's transcription); the dictionary glosses the last 'to come together'.
\end{cousinweb}

\begin{grammarlens}[title={What you've built}]
One more word for what I like.
\end{grammarlens}

\subsection*{Your turn}
\begin{itemize}
\raggedright
  \item \emph{Say it:} \emph{yīnyuèhuì}
  \item \emph{Say it:} \emph{yīnyuèhuì}, once more
  \item \emph{Say it:} \emph{míngtiān wǎnshang wǒmen qù yīnyuèhuì} — Tomorrow evening we're going to a concert.
  \item \emph{Recall:} say the Mandarin for song, then the Mandarin for pleasant to hear, then say \emph{yīnyuèhuì} again
\end{itemize}

\begin{tcolorbox}[breakable,colback=teal!4,colframe=teal!35!black,title={Before you move on}]
What does \zh{音乐会} mean? (\textquotedblleft{}Concert\textquotedblright{}.) Say it once more.
\end{tcolorbox}
\section[jítā]{jítā — guitar}

\label{lesson:ZH-C128-jita}



\begin{quote}
Before the new one: say the Mandarin for pleasant to hear, then the Mandarin for concert.
\end{quote}

\subsection*{You'll want to know: jítā}
\emph{jítā} — \textquotedblleft{}guitar\textquotedblright{}.

\begin{quote}
\emph{wǒ xiǎng xué jítā} — I'd like to learn the guitar.
\end{quote}

Its characters come later; for now it is a word for the ear.

\begin{cousinweb}
\emph{jítā} is usually traced to English \emph{guitar}; its two characters were chosen for their sound, the first a word for lucky and the second the \zh{他} of 'he'.
\end{cousinweb}

\begin{grammarlens}[title={What you've built}]
One more word for what I like.
\end{grammarlens}

\subsection*{Your turn}
\begin{itemize}
\raggedright
  \item \emph{Say it:} \emph{jítā}
  \item \emph{Say it:} \emph{jítā}, once more
  \item \emph{Say it:} \emph{wǒ xiǎng xué jítā} — I'd like to learn the guitar.
  \item \emph{Recall:} say the Mandarin for pleasant to hear, then the Mandarin for concert, then say \emph{jítā} again
\end{itemize}

\begin{tcolorbox}[breakable,colback=teal!4,colframe=teal!35!black,title={Before you move on}]
What does jítā mean? (\textquotedblleft{}Guitar\textquotedblright{}.) Say it once more.
\end{tcolorbox}
\section[(dialogue)]{Looking back: from What People Are Like to Music}

\label{lesson:ZH-R128-review-123-128}



\begin{quote}
Say the Mandarin for clever and conscientious.
\end{quote}

\subsection*{Your turn}
\begin{itemize}
\raggedright
  \item \emph{Say it:} clever, conscientious, lovable, warm-hearted, to care about
  \item \emph{Say it:} a person's height, hair, handsome, ugly, to resemble
  \item \emph{Say it:} to get angry, to worry, to be at ease, anxious, upset
  \item \emph{Say it:} often, sometimes, in the past, a moment ago, already
  \item \emph{Say it:} these, those, which ones, a few, like this
  \item \emph{Say it:} music, song, pleasant to hear, concert, guitar
  \item \emph{Read it:} \zh{以} \zh{些} \zh{音} — each one aloud, then a word that holds it
\end{itemize}

\begin{tcolorbox}[breakable,colback=teal!4,colframe=teal!35!black,title={Before you move on}]
Clever? (\emph{cōngming}.) Conscientious? (\emph{rènzhēn}.) Lovable? (\emph{kě'ài}.) Warm-hearted? (\emph{rèqíng}.) To care about? (\emph{guānxīn}.) A person's height? (\textbf{\zh{个子}}, \emph{gèzi}.) Hair? (\emph{tóufa}.) Handsome? (\emph{shuài}.) Ugly? (\emph{nánkàn}.) To resemble? (\emph{xiàng}.) To get angry? (\textbf{\zh{生气}}, \emph{shēngqì}.) To worry? (\emph{dānxīn}.) To be at ease? (\emph{fàngxīn}.) Anxious? (\emph{zháojí}.) Upset? (\emph{nánguò}.) Often? (\emph{chángcháng}.) Sometimes? (\emph{yǒushíhou}.) In the past? (\textbf{\zh{以前}}, \emph{yǐqián}.) A moment ago? (\emph{gāngcái}.) Already? (\emph{yǐjīng}.) These? (\textbf{\zh{这些}}, \emph{zhèxiē}.) Those? (\textbf{\zh{那些}}, \emph{nàxiē}.) Which ones? (\textbf{\zh{哪些}}, \emph{nǎxiē}.) A few? (\textbf{\zh{一些}}, \emph{yìxiē}.) Like this? (\emph{zhèyàng}.) Music? (\textbf{\zh{音乐}}, \emph{yīnyuè}.) Song? (\emph{gē}.) Pleasant to hear? (\emph{hǎotīng}.) Concert? (\textbf{\zh{音乐会}}, \emph{yīnyuèhuì}.) Guitar? (\emph{jítā}.)
\end{tcolorbox}
